\documentclass[11pt,a4paper]{article}
\usepackage{fshf-style}
\SetFshfShortTitle{Equities Portfolio Research}

\begin{document}

\makefshftitle{Equities Portfolio Research}%
{Equities --- United States and Europe}%
{Timur Khairullin, Anton Hilger \& Sandro Janashvili}%
{24 June 2026}%
{As of 24 June 2026 \textperiodcentered\ Holding period: deployment to the December 2026 rebalancing \textperiodcentered\ Coverage: United States and Europe}

\setcounter{tocdepth}{2}
\tableofcontents
\clearpage

\section*{Executive Summary}
\addcontentsline{toc}{section}{Executive Summary}

Halfway through 2026, the global economy is mostly driven by one dominant factor -- the geopolitical situation caused by the war in Iran. The most significant outcome to note is the supply side energy shock burdening key central banks at their different cyclical starting points. The shock pushed inflation above target everywhere simultaneously and forced a synchronized hawkish turn: the ECB delivered its first hike since 2023, lifting the deposit rate to 2.25\%, while the Federal Reserve held at 3.50--3.75\% and removed any bias toward cuts. Thus, the current environment negatively affects those companies, whose value sits in future profits, while rewarding cheap, cash generative, rate-benefiting, and pricing power sectors.

This paper first presents the full underlying research: the global macro and rate environment, the US sector analysis, and the European sector analysis. Each part first includes its complete data set and figures and, second, synthesizes them into a single combined allocation. We construct the equity portion of the portfolio, which accounts for 100 percent of the portfolio, around a passive core of ETFs, approximately 60\% of which consists of the broad U.S.\ and European markets, with the remaining approximately 40\% consisting of thematic satellites: European banks, global defense, health care, consumer staples/food, energy, US financials, and a small rare-earths position.

The portfolio's regional breakdown is slightly skewed toward the U.S.\ (\(\sim\)55\%) due to the growth momentum and liquidity of thematic sectors, which is offset by the discount in European asset valuations, the ECB's interest rate policy, which creates favorable conditions for banks, and the fact that the European assets carry no USD/EUR currency risk for the euro-denominated fund --- unlike the unhedged US sleeve, which carries a negative currency contribution from the dollar's 2026 weakness. The portfolio's largest active holdings (energy and the defense industry) represent, in part, a volatile forecast on the duration of the energy crisis. However, at the time of writing this paper, that crisis is gradually subsiding, although not without the emergence of new risks.

A US--Iran interim ceasefire (a memorandum of understanding) was signed on 17 June, and front-month Brent futures fell back toward \$80 by 19 June from a late March peak near \$119. Nevertheless, it is hard to claim if it was the final accord in this long-standing global uncertainty, especially given that de-escalation is not yet settled: by 20--21 June Iran again declared the Strait of Hormuz closed, while US disputing that fact. Hence, we size energy down and reframe it as a hedge, while retaining defense on its more durable NATO/budget backlog logic due to its unreliability on the previously discussed factors with financing being solely driven by the treaty's budget commitments.

\section{Methodology and Source Audit}

All macro and instrument figures were cross checked against primary or institutional open sources. Three corrections emerged and are reflected throughout:

\begin{itemize}
\item \textbf{Fed:} Even though 17 June FOMC held 3.50--3.75\%, its dot plot removed the bias toward cuts and pushed reductions into 2027--28; this was Kevin Warsh's first meeting as chair, and markets moved to price the next hike as early as October. (Federal Reserve; CNBC; Advisor Perspectives, June 2026.)
\item \textbf{Peace path:} A 60-day ceasefire between the US and Iran signed on June 17 intended to pause the conflict and reopen the Strait of Hormuz. Oil markets responded immediately, which could be seen in front-month Brent futures dropping by roughly 8.5\% in a week to around \$80, and Goldman Sachs cutting its Q4 price forecast to the same level. However, the physical spot price (as tracked by EIA/FRED) was still around \$93 on June 11, a level recorded eight days before the \(\sim\)\$80 futures print; the gap between the two therefore reflects the decline in prices over that week as the ceasefire firmed, rather than a lag between spot and futures markets on the same date. Nevertheless, the calm and greater confidence in how events would unfold did not last long: by June 20--21, Iran had once again announced the closure of the strait, and the U.S.\ responded by attempting to confound. Thus, a consensus was still lacking, and technical negotiations continued in Switzerland.
\item \textbf{ECB:} confirmed +25bp to a 2.25\% deposit rate effective 17 June, with 2026 growth cut to 0.8\% and a further July hike seen as likely. (ECB statement, 11 June 2026; Capital Economics.)
\end{itemize}

The valuation premise holds: Europe trades at a 9--26\% forward-P/E discount to the US with roughly double the dividend yield. For a euro-based fund the dollar's 2026 weakness against the euro considerably slows down the unhedged US exposure. (InvestingBrokers; eupersonalfinance.eu, 2026.)

\section{The Global Macro and Rate Environment}

The proximate cause of the 2026 inflation impulse is the Iran conflict and the disruption to oil and gas supply, including fears over the Strait of Hormuz. Brent began the year near \$66, climbed to a peak around \$119 at the end of March, then eased back as peace expectations built. It was followed by front-month futures reaching roughly \$80 by 19 June, although EIA/FRED was still around \$93 (\$92.84) on 11 June. This round-trip is the single most important macro development of the year: it lifted headline inflation everywhere, fed into core measures with a lag, and introduced two-sided risk, since any breakdown or consolidation of the peace process moves the global price level sharply either way.

\fshffig{equity-fig1-brent-2026.png}{Brent crude during the Middle East conflict (2026)}{Sources: EIA/FRED (spot); ICE/Investing.com (front-month futures).}

\subsection{Inflation above target across the major economies}

The energy shock left inflation uncomfortably above the 2\% objective everywhere. That brought about the US core PCE to rise from 3.0\% in December 2025 to roughly 3.3\% by April 2026. Euro-area headline HICP accelerated to 3.2\% in May, the highest since 2023, with core at 2.5\%. UK CPI ran around 3.3\% with services inflation stabilizing near 4.5\%. Long-run inflation expectations stayed comparatively well anchored, which can be comprehended as signal that investors read the spike as a supply-driven episode rather than a de-anchoring.

\fshffig{equity-fig2-inflation-above-target.png}{Inflation above target across major economies (2026)}{Sources: BEA, Eurostat, ONS.}

\subsection{Divergence among the major central banks}

The shared shock produced divergent responses. The Fed held its target range at 3.50--3.75\%, with the conversation tilting toward a possible hike. The ECB raised by 25bp on 11 June, lifting the deposit rate to 2.25\% --- its first increase since September 2023 and a reversal of eight consecutive cuts. The Bank of England held at 3.75\% with a hawkish split. The Bank of Japan, still the outlier on level, extended its gradual normalization with a 25bp hike to 1.00\% on 16 June (a 7--1 vote) --- its first increase since December and its highest policy rate since 1995. The euro area illustrates the dilemma most acutely: the bloc contracted slightly in Q1 2026, yet staff revised inflation up and growth down.

\fshffig{equity-fig3-policy-rates.png}{Major central bank policy rates (as of 20 June 2026)}{Sources: respective central banks.}

\subsection{US policy path and the term structure}

The dollar market anchors global financing conditions, which results in the US trajectory's outsized weight: having cut about 175 basis points across 2024--25, the Fed entered 2026 on hold, and the energy-driven reacceleration removed the basis for further easing. The Treasury curve normalized into a conventional upward slope, keeping the cost of term capital elevated worldwide.

\fshffig{equity-fig4-fed-funds-upper.png}{US federal funds target, upper bound (2024--2026)}{Source: FOMC statements.}

\subsection{Financial conditions and outlook}

Despite the hawkish backdrop, broad financial conditions remained comparatively benign: US high-yield credit spreads stayed compressed near 2.7 percentage points, well below crisis levels, and equities recovered on improving risk sentiment as de-escalation hopes grew. The central tension into H2 2026 is between still-easy financial conditions and a policy stance restrictive in intent. The dominant risk is two-sided and runs through energy: a durable settlement lets inflation resume falling and reopens the path to easing; renewed escalation around Hormuz could force central banks from holding into active tightening --- leaving the global rate environment in a higher-for-longer regime with an unusually wide range of outcomes.

\fshffig{equity-fig5-treasury-curve.png}{US Treasury yield curve (mid-June 2026)}{Sources: US Treasury constant-maturity; Fed H.15.}

\section{United States --- Equity Sector Research}

The US economy is in a late-cycle regime: firm, accelerating growth and rising, sticky inflation, with manufacturing at a multi-year high and headline CPI above 4\% on an oil-driven energy shock. The Fed holds at 3.50--3.75\%; futures had priced meaningful hike risk by December with no cuts in sight. This higher-for-longer environment favors cheap, cash-generative, and defensive sectors over expensive, rate-sensitive ones.

\subsection{Regime: policy, growth, inflation}

\begin{itemize}
\item \textbf{Policy} --- target range held at 3.50--3.75\%; the forward curve priced a \(\sim\)59\% probability of at least one hike by December (as of mid-June), with the priced probability of cuts near 1\%. The stance shifted from neutral hold to hawkish hold within a month.
\item \textbf{Growth} --- ISM Manufacturing PMI reached 54.0 in May, a fifth month above 50 consistent with \(\sim\)2.2\% annualised GDP; New Orders rose to 56.8. Payrolls added 172,000 vs \(\sim\)85,000 expected; unemployment held at 4.3\%.
\item \textbf{Inflation} --- headline CPI rose 4.2\% YoY in May, the fastest since 2023; energy was over 60\% of the monthly gain. Core PCE reached 3.3\% YoY in April, 130bp above target --- sticky, not solely energy-driven.
\end{itemize}

Regime classification: growth rising, inflation rising --- the reflation / late-cycle quadrant. It favors commodity-linked and cyclical sectors (Energy, Materials, Industrials) and defensives with pricing power (Staples, Health Care); it penalizes long-duration rate-sensitive sectors (Information Technology, Real Estate, Utilities as bond proxy).

\subsection{US key macro indicators}

\begin{table}[htbp]
\caption{US key macro indicators}
\small
\begin{tabularx}{\linewidth}{P{3.4cm} P{2.9cm} P{3.4cm} L}
\toprule
\tabh{Indicator} & \tabh{Reading} & \tabh{Signal} & \tabh{Implication} \\
\midrule
Fed Funds Rate & 3.50--3.75\% & Hold (97.4\% prob, Jun 17) & Higher-for-longer; hike risk by Dec \\
Dec 2026 hike probability & \(\sim\)59\% & Hawkish drift & No cuts priced; possible tightening \\
10-year Treasury yield & \(\sim\)4.54\% & Elevated & Headwind for rate-sensitive sectors \\
2s10s curve slope & +0.38\% & Positive / steepening & Tailwind for bank NIM \\
US Dollar (DXY) & \(\sim\)99.7 & Countertrend bounce (+1.8\% MoM) within 2026 weakness vs EUR & Headwind for commodities, Materials \\
ISM Manufacturing PMI & 54.0 (May) & Expanding (5th month) & GDP growth \(\sim\)2.2\% annualised \\
New Orders sub-index & 56.8 & Rising from 54.1 & Production momentum continues \\
Non-farm payrolls (May) & +172,000 & Beat vs \(\sim\)85,000 exp. & Labor market firm \\
Unemployment rate & 4.3\% & Stable & No slack emerging \\
CPI (May, YoY) & +4.2\% & Fastest since 2023 & Inflation re-accelerating \\
CPI energy contribution & \(>\)60\% of gain & Oil shock (Iran) & Geopolitical risk in headline \\
Core PCE (Apr, YoY) & +3.3\% & Highest since late 2023 & Underlying inflation sticky \\
Headline PCE (Apr, YoY) & +3.8\% & 2nd consecutive accel. & Broad price pressure \\
\bottomrule
\end{tabularx}
\fshfsource{Sources: CME FedWatch; FRED; BLS; BEA; ISM. Data as of 13--16 June 2026.}
\end{table}

\subsection{US sector analysis}

Each sector is assessed on valuation vs the \(\sim\)21\(\times\) market multiple, earnings-revision direction, and regime fit. Calls are relative to the S\&P 500. The valuation spread below is the starting point: Energy and Financials screen cheapest, while Information Technology and Consumer Discretionary carry the richest multiples and the greatest rate-sensitivity. These sector convictions are screened on US sector data but are implemented, per the PRIIPs constraint documented in Section~7, through UCITS vehicles that are predominantly European sector ETFs; the European sector analysis in Section~4.4 carries the corresponding convictions for the vehicles that ship, and only the US financials sleeve retains direct US sector exposure.

\begin{table}[htbp]
\caption{US sector assessment relative to the S\&P 500}
\small
\begin{tabularx}{\linewidth}{P{2.6cm} P{0.9cm} P{2.9cm} P{3.5cm} L}
\toprule
\tabh{Sector (ETF)} & \tabh{Call} & \tabh{Valuation vs market} & \tabh{Earnings / regime} & \tabh{Key risk} \\
\midrule
Energy (XLE) & OW & P/E 12.5\(\times\) --- cheapest & Rising with oil; favorable & Ceasefire collapses oil premium \\
Industrials (XLI) & OW & P/E 27\(\times\) --- expensive & Backlog \(>\)1 BTB; favorable & ISM rollover; budget delay \\
Health Care (XLV) & OW & P/E 18\(\times\) --- below market & Demographic; favorable & Drug pricing legislation \\
Consumer Staples (XLP) & OW & P/E 20\(\times\) --- at market & Food pricing power & Disinflation; GLP-1 drag \\
Materials (XLB) & OW\textsuperscript{\dag} & P/E 18\(\times\) --- below market & Mixed; dollar headwind & Stronger dollar \\
Financials (XLF) & N/OW & P/E 15.5\(\times\) --- below market & NIM tailwind; favorable & Credit deterioration \\
Comm.\ Services (XLC) & N & P/E 16\(\times\) --- below market & Ad growth solid; neutral & Ad-spend cyclicality \\
Utilities (XLU) & N & P/E 18\(\times\) --- below market & AI power demand; balanced & Rising 10-year yield \\
Info.\ Technology (XLK) & N/UW & P/E 40\(\times\) --- most expensive & Strong EPS; unfavorable & Hawkish Fed compresses multiple \\
Consumer Discr.\ (XLY) & UW & P/E 30\(\times\) --- expensive & Worsening; unfavorable & Resilient-consumer surprise \\
Real Estate (XLRE) & UW & Yld 3.4\% vs 4.5\% 10y & Pressured; unfavorable & Fed pivot toward cuts \\
\bottomrule
\end{tabularx}
\fshfsource{Sources: stockanalysis.com; S\&P Global; FactSet. As of 16 June 2026.}
\end{table}

\fshffig{equity-figA-us-forward-pe.png}{US sector forward P/E vs the market multiple}{Sources: stockanalysis.com; S\&P Global; FactSet.}

\subsection{US conviction ranking and portfolio weights}

\begin{table}[htbp]
\caption{US conviction ranking and active tilts}
\small
\begin{tabularx}{\linewidth}{P{0.5cm} P{3.1cm} L Q{2.6cm}}
\toprule
\tabh{\#} & \tabh{Sector} & \tabh{Call \& rationale} & \tabh{Active tilt vs S\&P 500} \\
\midrule
1 & Energy (XLE) & OW --- cheapest sector; all lenses aligned; inflation hedge & +2.0 \\
2 & Industrials (XLI) & OW --- defense backlog; NATO structural tailwind & +2.0 \\
3 & Health Care (XLV) & OW --- best risk-adjusted: cheap, low rate-sensitivity, demographic & +2.5 \\
4 & Consumer Staples (XLP) & OW --- defensive anchor; food pricing power & +2.0 \\
5 & Materials (XLB) & OW --- rare-earths policy tailwind; dollar headwind limits to names & +1.0 \\
6 & Financials (XLF) & N/OW --- curve re-steepening supports NIM; credit risk caps conviction & +1.0 \\
7 & Comm.\ Services (XLC) & N --- lower-valuation tech proxy; advertising resilient & 0.0 \\
8 & Utilities (XLU) & N --- rate headwind balanced by AI power demand & 0.0 \\
9 & Info.\ Technology (XLK) & N/UW --- best EPS growth but richest multiple, highest rate-sensitivity & \(-\)6.0 \\
10 & Consumer Discr.\ (XLY) & UW --- expensive; consumer squeezed by inflation + borrowing costs & \(-\)1.5 \\
11 & Real Estate (XLRE) & UW --- most direct rate loser; yield below 10-year Treasury & \(-\)1.5 \\
\bottomrule
\end{tabularx}
\fshfsource{Source: S\&P Dow Jones Indices, late May 2026.}
\end{table}

\fshffig{equity-figB-us-active-tilts.png}{US sector active tilts --- overweights vs underweights}{\relax}

\section{Europe --- Equity Sector Research}

European equities trade at a historic valuation discount to the US, a compelling setup for dividend-yield-oriented exposure; the STOXX 600 vs S\&P 500 forward-earnings discount stood near 20\%. On 11 June the ECB raised its policy rate 25bp to 2.25\% in reaction to the energy-driven landscape, a restrictive shift that improves bank margins while penalizing rate-sensitive sectors. The European sleeve is built from passive funds and ETFs with targeted sector overweights in Banking, Energy, and Defense.

\subsection{Growth, fund flows, and sector rotation}

Major European indices showed solid twelve-month growth (STOXX Europe 600 +18.6\%, STOXX Europe 50 +20.0\%). Fund flows underscored the momentum: Q1 2026 volume rose 19.4\% QoQ (EUR 29.9bn) to EUR 184.2bn total, with equity funds drawing EUR 95.3bn (up from EUR 42.9bn in Q4 2025); passive funds and ETFs captured 65\% of total flows. Rising AI skepticism drove EUR 23.3bn out of growth-focused funds in Q1 --- capital rotating from Growth toward defensive, value, and dividend-yielding European assets.

\fshffig{equity-fig6-european-fund-flows.png}{European fund flows: rotation into passive equity (2026)}{Source: Morningstar Direct, data to 31 Mar 2026.}

\subsection{European key macro indicators}

\begin{table}[htbp]
\caption{European key macro indicators}
\small
\begin{tabularx}{\linewidth}{P{3.2cm} P{3.1cm} P{3.2cm} L}
\toprule
\tabh{Indicator} & \tabh{Reading (2026)} & \tabh{Signal} & \tabh{Implication} \\
\midrule
ECB Policy Rate & 2.25\% & 25bp hike (Jun 11) & Improved bank margins; RE headwind \\
STOXX Europe 600 & +10.14\% (6m) & Solid growth & Positive European momentum \\
Valuation Discount & 19.98\% & STOXX 600 vs S\&P 500 & Attractive for dividend-yield \\
Equity Fund Inflows & EUR 95.3bn & +122\% vs Q4 2025 & Strong interest; passive dominance \\
Growth Equity Funds & \(-\)EUR 23.3bn (Q1) & Rising AI skepticism & Rotation away from Tech/Growth \\
\bottomrule
\end{tabularx}
\fshfsource{Sources: Morningstar; ecb.europa.eu.}
\end{table}

\subsection{European fundamental valuation}

\begin{table}[htbp]
\caption{European fundamental valuation vs the US}
\small
\begin{tabularx}{\linewidth}{P{2.9cm} P{3.0cm} P{2.6cm} L}
\toprule
\tabh{Metric} & \tabh{STOXX Europe 600} & \tabh{S\&P 500 (US)} & \tabh{Market implication} \\
\midrule
Forward P/E & 18.02\(\times\) & 22.52\(\times\) & \(\sim\)20\% discount on expected 12m earnings; attractive value entry \\
Dividend Yield & 2.49\% & 1.05\% & More than double the US payout; downside protection in volatile regimes \\
\bottomrule
\end{tabularx}
\fshfsource{Sources: Reuters; macromicro.me; justETF.}
\end{table}

\subsection{European sector calls and conviction}

\begin{table}[htbp]
\caption{European sector calls}
\small
\begin{tabularx}{\linewidth}{P{2.9cm} P{0.9cm} L P{4.0cm}}
\toprule
\tabh{Sector} & \tabh{Call} & \tabh{Rationale \& regime fit} & \tabh{Key risk} \\
\midrule
Financials (Value) & OW & Direct beneficiary of ECB hike to 2.25\%; improves bank NIM --- favorable & Sudden cuts in a recession \\
Energy \& Defense & OW & Geopolitical uncertainty (US--Iran, Russia--Ukraine) favors contractors --- favorable & Framework peace collapses risk premium \\
Real Estate & UW & Highly sensitive to financing costs; 2.25\% weighs on valuations --- unfavorable & Unexpectedly fast dovish ECB pivot \\
Consumer Discr. & UW & Higher rates lower demand, pressuring spending --- unfavorable & Rapid recovery in domestic demand \\
\bottomrule
\end{tabularx}
\fshfsource{Source: ecb.europa.eu.}
\end{table}

\subsection{European source-sleeve portfolio weights}

\begin{table}[htbp]
\caption{European source-sleeve portfolio weights}
\small
\begin{tabularx}{\linewidth}{L P{1.4cm} P{3.1cm} Q{1.4cm}}
\toprule
\tabh{Sector / ETF} & \tabh{Call} & \tabh{Focus} & \tabh{Weight} \\
\midrule
iShares STOXX Europe 600 UCITS ETF & Core & Broad market & 40\% \\
iShares STOXX Europe 50 UCITS ETF & Core & Blue chips & 20\% \\
iShares STOXX Europe 600 Banks & OW & Financials / Value & 20\% \\
iShares MSCI Europe Energy Sector & OW & Energy & 10\% \\
iShares Europe Defence UCITS ETF & OW & Defense & 10\% \\
\bottomrule
\end{tabularx}
\fshfsource{Source: regional research note.}
\end{table}

\fshffig{equity-fig7-source-european-sleeve.png}{Source European sleeve allocation (regional note)}{\relax}

This source-sleeve table reproduces the regional research note; where it differs from the combined portfolio in Sections 5 and 6, which uses EXSA for the Europe core, EXH1 for energy, and the HANetf NATO fund for defense, and carries no STOXX Europe 50 position, the Section~6 list is the one that governs deployment.

\section{Synthesis: Combined Regime and Thesis}

It is worth observing at the outset that the three analyses, however distinct in their framing, ultimately describe a single phenomenon: one global energy shock passing through three central banks at different points in the economic cycle. The United States finds itself in the rising-growth, rising-inflation quadrant typical of late-cycle reflation. The euro area, by contrast, occupies the falling-growth, rising-inflation territory that characterises mild stagflation. Hence, the shock itself is identical, but the cyclical positions are not, and that divergence has material consequences for sector selection.

In the American context, portfolio construction can afford to ride the cycle, although In Europe, the priority shifts to sorting those sectors that benefit from higher rates from those that suffer under them. Yet for all these regional differences, the two analyses converge on the same recommendations where the evidence is strongest. Banks and financials merit an overweight in both markets. In addition, real estate and expensive, long-duration technology merit an underweight in both.
\[ \text{Stance} = f\,(\, \text{starting rate},\ \text{growth},\ \text{credibility risk}\,) \]

\subsection{Combined cross-regional conviction}

\begin{table}[htbp]
\caption{Combined cross-regional conviction}
\small
\begin{tabularx}{\linewidth}{P{2.7cm} P{2.9cm} L P{2.5cm}}
\toprule
\tabh{Theme} & \tabh{Expression} & \tabh{Combined conviction} & \tabh{Status (June audit)} \\
\midrule
Banks / Financials & EU banks + US financials & Highest --- NIM tailwind both regions & Strengthened \\
Defense & Global NATO-aligned ETF & High --- budget backlog, not oil-dependent & Mostly intact \\
Health Care & Cheap defensive & Medium-high --- best risk-adjusted defensive & Unchanged \\
Staples / Food & Pricing power & Medium & Unchanged \\
Energy & Oil \& gas & Medium \(\rightarrow\) reduced & Weakened \\
Rare Earths & Strategic-metals ETF & Low-medium --- China-concentration caveat & Unchanged \\
Technology & (core only) & Underweight --- richest multiple & Reinforced \\
Real Estate & (core only) & Underweight --- rate-penalized both regions & Reinforced \\
\bottomrule
\end{tabularx}
\end{table}

\textbf{The energy/defense decoupling.} The source notes treated energy and defense as one geopolitical block, which have now decoupled. Defense rests on NATO's path toward 5\% of GDP by 2035 and multi-year US budget backlog --- visibility that survives a falling oil price. Energy rests directly on the oil premium, much of which unwound through mid-June as front-month futures fell toward \$80 --- though the renewed Hormuz dispute on 20--21 June shows how quickly that premium can rebuild. We keep defense as a conviction overweight and demote energy to a smaller, tactical hedge against the ceasefire collapsing.

\textbf{Rare-earths caveat.} The liquid UCITS vehicle (VanEck) is China/Australia/US-heavy, so it imperfectly expresses the Western de-risking thesis; sized small (3\%) with that mismatch flagged.

\subsection{Regional allocation}
\[ \text{US weight} \propto \frac{\text{growth momentum} + \text{theme liquidity}}{\text{valuation} + \text{FX drag}} \]

The US pulls weight through growth momentum and deeper theme liquidity; Europe through the \(\sim\)15--20\% valuation discount, double the dividend yield, the ECB-rate tailwind to its large bank sector, and no FX drag for a euro fund. With the energy theme softened and FX cutting against unhedged US exposure, an overwhelming US tilt is weaker than a US-only read implies. We set \(\sim\)55\% US / 45\% Europe --- a deliberate, moderate tilt. US positions are held unhedged as a known, accepted exposure.

\section{Portfolio Construction}

A passive broad-market core captures the valuation discount and market return at minimal cost; sector/thematic ETF satellites express the macro convictions while avoiding single-stock idiosyncratic risk.

\begin{table}[htbp]
\caption{Recommended combined portfolio --- 60\% core / 40\% satellites}
\small
\begin{tabularx}{\linewidth}{P{1.6cm} P{4.6cm} Q{1.3cm} L}
\toprule
\tabh{Sleeve} & \tabh{Position} & \tabh{Weight} & \tabh{Rationale} \\
\midrule
Core & US broad market (S\&P 500) & 34\% & Growth momentum; deepest, cheapest beta \\
Core & Europe broad market (STOXX 600) & 26\% & Valuation discount; income; no FX drag \\
Satellite & European banks & 8\% & Highest-conviction; ECB-hike NIM tailwind \\
Satellite & Global defense & 8\% & Structural NATO/budget backlog \\
Satellite & Health care & 6\% & Cheap, low rate-sensitivity defensive \\
Satellite & Consumer staples / food & 6\% & Pricing power vs sticky inflation \\
Satellite & Energy (oil \& gas) & 5\% & Tactical hedge vs ceasefire collapse \\
Satellite & US financials & 4\% & Curve steepening supports margins \\
Satellite & Rare earths / strategic metals & 3\% & Small thematic tilt; China caveat \\
\bottomrule
\end{tabularx}
\end{table}

Banks-plus-financials (12\% combined) is the largest active bet. Technology and real estate carry no satellite weight, held below benchmark via the core --- the funding source for the overweights. Blended trailing dividend yield is 1.9\%, roughly double the 1.0\% of the S\&P 500; portfolio beta is 0.86 against MSCI World in EUR and 0.73 against the S\&P 500 in EUR --- market-like exposure with the intended value/defensive tilt.\footnote{Computed 24 July 2026 from Yahoo Finance listings of the Section 7 instruments at the table weights. Yield: weight $\times$ per-ETF trailing-12-month cash distributions over the last close, taken from the distributing German listings where available and from distributing same-index proxies (SPY, XLF, REMX) for accumulating share classes; the defense sleeve has no distributing twin, so the blend covers 92\% of portfolio weight, renormalized. Beta: OLS of the weighted portfolio weekly EUR return series (Xetra listings, 104 weeks to 24 July 2026) on iShares Core MSCI World (EUNL.DE) and on the EUR-listed S\&P 500 tracker (SXR8.DE); converting the S\&P 500 index at spot EUR/USD instead gives 0.62, the difference reflecting asynchronous closing times.}
\[ \mathrm{E}[R_p] = \sum_i w_i \cdot \mathrm{E}[R_i]\,, \qquad \text{with } \sum_i w_i = 1 \]

\fshffig{equity-fig8-combined-portfolio.png}{Recommended combined portfolio --- 60\% core / 40\% satellites}{\relax}

\section{Instrument List (IBKR-Deployable UCITS Vehicles)}

Critical constraint: under EU PRIIPs/KID rules, US-domiciled ETFs (XLE, XLK, ITA, REMX-US) are generally not purchasable by EU-resident clients on Interactive Brokers. The list uses UCITS equivalents. All rows were checked against fund-provider and exchange data (ISINs verified against issuer product pages, July 2026); listings should still be re-confirmed in IBKR on deployment day.

\begin{table}[htbp]
\caption{Instrument list --- IBKR-deployable UCITS vehicles}
\small
\begin{tabularx}{\linewidth}{P{1.9cm} L P{1.2cm} P{2.6cm} P{1.1cm} P{1.3cm}}
\toprule
\tabh{Sleeve} & \tabh{Fund} & \tabh{Ticker} & \tabh{ISIN} & \tabh{TER} & \tabh{Status} \\
\midrule
US core & iShares Core S\&P 500 UCITS ETF (Acc) & CSPX & IE00B5BMR087 & 0.07\% & verified \\
Europe core & iShares STOXX Europe 600 UCITS ETF (DE) & EXSA & DE0002635307 & 0.20\% & verified \\
EU banks & iShares STOXX Europe 600 Banks UCITS ETF (DE) & EXV1 & DE000A0F5UJ7 & 0.47\% & verified \\
Defense & HANetf Future of Defence UCITS ETF (Acc) & NATO & IE000OJ5TQP4 & 0.49\% & verified \\
Rare earths & VanEck Rare Earth \& Strategic Metals UCITS & VVMX & IE0002PG6CA6 & 0.59\% & verified \\
Health care & iShares STOXX Europe 600 Health Care UCITS (DE) & EXV4 & DE000A0Q4R36 & 0.46\% & verified \\
Energy & iShares STOXX Europe 600 Oil \& Gas UCITS (DE) & EXH1 & DE000A0H08M3 & 0.47\% & verified \\
Staples / food & iShares STOXX Europe 600 Food \& Beverage (DE) & EXH3 & DE000A0H08H3 & 0.46\% & verified \\
US financials & iShares S\&P 500 Financials Sector UCITS ETF & IUFS & IE00B4JNQZ49 & 0.15\% & verified \\
\bottomrule
\end{tabularx}
\end{table}

Liquidity / cost: the verified core/bank/defense lines are large, liquid funds with tight spreads; the sector-sleeve and rare-earth lines are thinner --- use limit orders and check spreads before sizing. Choose accumulating vs distributing share classes per the fund's tax preference (German- vs Irish-domiciled funds differ on withholding).

\section{Risk Assessment}

\subsection{The central scenario: the peace path}

The portfolio's largest active exposures make it, in net effect, a bet against a durable, rapid peace settlement --- a conscious, stated position, not an accident of sector selection.

\textbf{Scenario A --- ceasefire holds} (base case, but less certain after the 20--21 June Hormuz dispute). Oil settles into the \$70s, inflation resumes its decline, the path reopens to easing. Energy and (less so) defense premiums compress; rate-sensitive underweights reverse and drag. This is what the portfolio is most exposed to. Mitigants: energy already cut to 5\% and framed as a hedge; defense rests on backlog; the 60\% core participates in any tech/RE recovery.

\textbf{Scenario B --- ceasefire collapses, Hormuz re-escalation.} Oil spikes back, inflation re-accelerates, central banks move from holding into tightening. Energy and defense rip, banks benefit from higher-for-longer, underweights vindicated --- the portfolio outperforms.

Directional sensitivity: with energy \(\sim\)5\% and underweights via the core, a sustained \$10--15 Brent move is meaningful but not dominant --- sized so the portfolio tilts toward Scenario B rather than betting the book on it.

\subsection{Secondary risks}

\begin{itemize}
\item \textbf{Hawkish Fed surprise} --- supports financials but pressures the core; offset by the value/defensive tilt.
\item \textbf{Deeper European recession} --- the genuine stagflation risk; hurts EU banks and the EU core; cushioned by the discount and dividend.
\item \textbf{Defense budget delay} --- backlog is multi-year but order timing can slip; sized at 8\% to avoid single-point failure.
\item \textbf{Dollar moves} --- dollar strength hurts USD-priced rare-earths/commodity exposure but lifts the EUR value of the unhedged US sleeve; dollar weakness does the opposite, dragging on unhedged US returns in EUR --- a two-sided exposure.
\item \textbf{Rare-earths China concentration} --- a China shock hits the vehicle; mitigated by the small 3\% size.
\end{itemize}

\subsection{What would make us wrong}

A permanent, verified peace agreement (not just the 60-day interim), Hormuz fully reopened, and oil stabilizing in the \$70s or lower. In that world the regime thesis weakens, and the portfolio should rotate toward the rate-sensitive sectors we currently underweight. The December rebalancing --- or an earlier trigger --- is the decisive point.

\section{Monitoring Plan and Conclusion}

Weekly watch items: Brent and US--Iran talks status; ECB/Fed communications and hike-expectation shifts (CME FedWatch, ECB-Watch); EUR/USD as a return driver; defense-budget news; satellite-ETF spreads/liquidity before resizing.

Pre-deployment checklist: (1) re-pull live oil price and ceasefire status on deployment morning; (2) re-confirm the instrument listings in IBKR; (3) refresh live forward-P/E and dividend figures and rescale weights to exactly 100\%; (4) confirm accumulating vs distributing share-class preference.

To summarize the position: the portfolio holds a value- and defensive-tilted global equity sleeve, supported by income and built around a passive core. The satellites express a coherent late-cycle, higher-for-longer thesis through targeted, liquid UCITS instruments. Their sizing reflects a deliberate choice to lean into the regime without concentrating risk on a single outcome. That restraint is warranted. The defining variable of mid-2026 is still unwinding, and the balance of probabilities could shift materially in either direction before December.

\section*{Limitations}
\addcontentsline{toc}{section}{Limitations}

Macro figures are verified to the date of writing, but oil and the ceasefire move daily --- re-check both at deployment. All instrument ISINs are verified against issuer data; re-confirm listings in IBKR at deployment. Valuation and yield figures should be refreshed to live data on the start date rather than carried from the June-11 source snapshots.

\begin{fshfreferences*}{Key Sources}
\addcontentsline{toc}{section}{Key Sources}
\item \textbf{Central Bank Policy}
\item European Central Bank --- Policy decisions and staff projections, 11 June 2026. \url{https://www.ecb.europa.eu}
\item Federal Reserve --- FOMC statement and Summary of Economic Projections, 17 June 2026. \url{https://www.federalreserve.gov}
\item Bank of England --- Monetary Policy Committee communications, June 2026. \url{https://www.bankofengland.co.uk}
\item Bank of Japan --- Statement on Monetary Policy, 16 June 2026. \url{https://www.boj.or.jp/en}
\item \textbf{Energy \& Commodity Prices}
\item EIA / FRED --- Brent crude spot price, Federal Reserve Bank of St.\ Louis. \url{https://fred.stlouisfed.org}
\item ICE / Investing.com --- Front-month futures data. \url{https://www.investing.com}
\item \textbf{Market \& Instrument Data}
\item Trading Economics --- Macroeconomic indicators and market data. \url{https://tradingeconomics.com}
\item Barchart --- Commodity and equity price data. \url{https://www.barchart.com}
\item Reuters --- Financial news and market data. \url{https://www.reuters.com}
\item Fortune --- Business and financial news. \url{https://fortune.com}
\item CNBC --- Financial news and market commentary. \url{https://www.cnbc.com}
\item Fox Business --- Financial news and market commentary. \url{https://www.foxbusiness.com}
\item Advisor Perspectives --- Investment research and market analysis. \url{https://www.advisorperspectives.com}
\item Capital Economics --- Independent macroeconomic research. \url{https://www.capitaleconomics.com}
\item InvestingBrokers --- Broker and instrument data. \url{https://www.investingbrokers.com}
\item eupersonalfinance.eu --- European retail investment and ETF data. \url{https://www.eupersonalfinance.eu}
\item justETF --- European ETF screener and fund data. \url{https://www.justetf.com}
\item Cbonds --- Fixed income and bond market data. \url{https://cbonds.com}
\item London Stock Exchange --- Exchange listings and instrument data. \url{https://www.londonstockexchange.com}
\item iShares / BlackRock --- ETF product data and fund documentation. \url{https://www.ishares.com}
\item HANetf --- Thematic and specialist ETF data. \url{https://www.hanetf.com}
\item VanEck --- ETF product data and fund documentation. \url{https://www.vaneck.com}
\end{fshfreferences*}

\fshfdisclaimer{on 24 June 2026}

\end{document}
